\sourceimport{finals-math-extra}{Prelude.h}
\section{Divisors and prime exponents}
\textbf{Choose the method first.} One moderate positive integer can use
trial division, removing each full prime power in $O(\sqrt n)$. Many
queries bounded by a feasible $V$ favor an SPF sieve; its memory is
$O(V)$. Trial division is a poor choice for many large inputs or
adversarial semiprimes. Neither method is unrestricted 64-bit
factorization.

\subsection{Generate divisors}
\texttt{positive\_divisors(n)} returns every positive divisor in sorted
order. It scans $d\le n/d$ and pairs $d$ with $n/d$, so it does not
construct divisors from a factorization. Runtime is $O(\sqrt n)$ plus
output storage. Requires $n>0$.
\sourceimport{finals-math-extra}{Divisors.h}

\subsection{Reason with exponents}
If $n=\prod p^{e_p}$, the least multiplier making $n$ a $k$th power is
$\prod p^{(k-e_p\bmod k)\bmod k}$. For $72=2^3\cdot3^2$, this gives
$2$ for squares. Pass the factorization from a suitable method; the
helper returns no value if the multiplier overflows signed 64-bit.
For factorial divisibility, Legendre's sum
$v_p(n!)=\lfloor n/p\rfloor+\lfloor n/p^2\rfloor+\cdots$
avoids constructing $n!$. Its helper requires prime $p$.
\sourceimport{finals-math-extra}{PrimeExponents.h}

\section{GCD classes and circular motion}
\subsection{Count exact pairwise GCDs}
Let $M_d$ count input values divisible by $d$. Every unordered pair
among those $M_d$ values has gcd a multiple of $d$. Thus, descending
from $V$, subtract counts already assigned to proper multiples:
$E_d=\binom{M_d}{2}-\sum_{j\ge2}E_{jd}$. Then $E_d$ counts pairs with
gcd exactly $d$, and $E_1$ is the coprime-pair count. This handles
duplicate values by index and excludes pairing an item with itself.
Requires positive values at most $V$; time $O(V\log V+N)$ and memory
$O(V)$. For a small number of pairs or a huge value range, use direct
pairwise gcd in $O(N^2\log V)$ instead. For $N<2$, every answer is 0.
\sourceimport{finals-math-extra}{GCDPairs.h}

\subsection{Are all pairwise GCDs equal?}
The supplied pair loop is correct for $N\ge2$, but costs
$O(N^2\log V)$ and reads out of bounds when $N<2$. Decide that boundary
from the statement: with no pairs, ``all pairs agree'' is vacuously true,
although there is no unique pairwise gcd value. A safe small-input version is:
\begin{lstlisting}
// Positive values. O(n^2 log(max a)).
bool same_gcd(const vector<long long>& a) {
    if (a.size() < 2) return true;
    long long g = std::gcd(a[0], a[1]);
    for (int i = 0; i < (int)a.size(); ++i)
        for (int j = i + 1; j < (int)a.size(); ++j)
            if (std::gcd(a[i], a[j]) != g)
                return false;
    return true;
}
\end{lstlisting}
For bounded positive values, \texttt{same\_pair\_gcd(a,V)} uses the
histogram. If every pair agrees, that common value must equal the gcd
of the whole array; check whether its exact class contains all
$N(N-1)/2$ pairs. Computing the whole-array gcd alone is insufficient:
$[6,10,15]$ has global gcd 1 but pair gcds 2,3,5. On $[6,10,14]$ every
pair gcd is 2; on $[2,2,4]$ it is also 2, so duplicates need not fail.

An alternative when SPF is already built: divide each number by the global
gcd, factor it, and reject if a prime appears in two different numbers.
The normalized numbers must be pairwise coprime. Mark each distinct prime
once per number; normalized 1 has no primes and may repeat. This is
$O(V+N\log V)$ with the existing linear SPF sieve, but the histogram is
more useful when the problem asks for counts of several gcd values.

\subsection{Orbits under repeated addition}
On $n$ circular positions, repeatedly adding $s$ forms
$\gcd(n,s)$ disjoint cycles, each of length $n/\gcd(n,s)$. For
$(n,s)=(12,8)$, there are four cycles of length three; from zero the
visited positions are $0,8,4$. Reduce reachability to the residue class
modulo the gcd, then check any safety condition along that cycle.
\sourceimport{finals-math-extra}{GCDOrbits.h}

\subsection{Base representation becomes divisibility}
If $N$ has exactly two equal digits $d$ in base $b$, then
$N=db+d=d(b+1)$. Enumerate positive divisors $d$ of $N$, recover
$b=N/d-1$, and require $b\ge2$ and $1\le d<b$. For $N=12$ this gives
$(d,b)=(1,11),(2,5),(3,3)$, but the last pair is illegal because a digit
must be smaller than its base. Thus the valid two-digit bases are 11 and 5.
This handles exactly two equal digits, not all palindromes. For the broader
2023 F task, smaller bases can have more digits and need direct conversion
and a palindrome check; partition the base ranges so results are not duplicated.
The proof is the positional numeral equation, not a guess from examples.
\href{https://icpcarchive.github.io/Latin%20American%20Contests/Latin%20American%20Regional%20Contest/2023%20Latin%20American%20Regional%20Contest/solution.pdf}{2023 official solutions}.

\subsection{Practice and adapt}
\href{https://cses.fi/problemset/task/1713}{Counting Divisors}: compare a
per-query factorization with an SPF or divisor-count sieve at the stated bounds.
\href{https://cses.fi/problemset/task/2417}{Counting Coprime Pairs}: use
$E_1$, not the number of prime inputs. Test repeated ones and repeated primes.
\href{https://cses.fi/problemset/task/2182}{Divisor Analysis}: the input
already supplies prime exponents; do not reconstruct or factor the huge integer.
Its modular divisor-product part requires additional exponent arithmetic.
These are practice tasks, not claimed new online-judge acceptances.
